\documentclass[10pt,a4paper]{article}
\usepackage[margin=22mm]{geometry}
\usepackage[T1]{fontenc}
\usepackage{lmodern}
\ifdefined\pdfmapfile\pdfmapfile{+lm.map}\fi
\usepackage{amsmath,amssymb,amsthm}
\usepackage{booktabs,graphicx}
\usepackage[expansion=false]{microtype}
\usepackage[numbers,sort&compress]{natbib}
\usepackage[hidelinks]{hyperref}
\renewcommand{\bibfont}{\small}
\setlength{\bibsep}{2pt}
\setlength{\parindent}{1.1em}
\setlength{\parskip}{0.13em}
\newtheorem{proposition}{Proposition}
\newcommand{\norm}[1]{\left\lVert #1\right\rVert}
\newcommand{\dd}{\mathrm{d}}
\newcommand{\q}{\mathbf q}
\newcommand{\vct}{\mathbf v}
\newcommand{\acc}{\mathbf a}
\newcommand{\force}{\mathbf F}
% Generated from the validated summary by benchmarks/hht/paper.py.
\graphicspath{{\detokenize{../results/hht/figures/}}}
\newcommand{\TrainMeshTime}{36.78}
\newcommand{\TrainCadTime}{1.077}
\newcommand{\TrainGain}{34.1}
\newcommand{\CamFineGain}{29.9}
\newcommand{\OscillatorOrders}{1.995, 2.000, 2.000}
\newcommand{\DropError}{0.0084}
\newcommand{\DropFineError}{0.0011}
\newcommand{\CamCoarseGapError}{3.536}
\newcommand{\CamFineGapError}{0.0540}
\newcommand{\TrainMatchError}{\ensuremath{7.49\times10^{-4}}}
\newcommand{\FineCamMeshError}{0.0842}
\newcommand{\CamCadTimeError}{0.0818}
\newcommand{\GuardOff}{7,202}
\newcommand{\GuardOn}{3,602}
\newcommand{\GuardGain}{1.48}
\newcommand{\StressCamGuard}{8.120}
\newcommand{\StressTrainGuard}{2.897}
\newcommand{\TrimFinestError}{\ensuremath{2.17\times10^{-5}}}
\newcommand{\TrimFinestMesh}{3.084}
\newcommand{\TrimCadTime}{46.883}
\newcommand{\CamTemporalFine}{0.0260}
\newcommand{\TrainEarlyCoarse}{0.312}
\newcommand{\TrainEarlyFine}{0.112}
\newcommand{\TrainLongCoarse}{1.508}
\newcommand{\TrainLongFine}{2.719}
\newcommand{\ProductionOraclePoses}{68}
\newcommand{\ProductionOracleAgreement}{68}
\newcommand{\EngineFigurePath}{\detokenize{../results/engine-comparison/engine-comparison.pdf}}
\newcommand{\EngineMuJoCoVersion}{4.0.1}
\newcommand{\EngineDrakeVersion}{1.58.0}
\newcommand{\EngineForceResidual}{\ensuremath{6.42\times10^{-13}}}
\newcommand{\EngineBenchmarkRows}{%
  STEP/HHT & 4.742 & 6.392 & 0.011708 & 0.013894 \\
  MuJoCo implicit & 10.955 & 12.345 & 0.007723 & 0.009202 \\
  MuJoCo RK4 & 1.388 & 1.478 & 0.000005 & 0.000011 \\
  Drake RK3 & 1.997 & 2.563 & 0.000300 & 0.000511 \\
}
\newcommand{\EngineTimingRows}{%
  STEP/HHT & 666.2 & 862.0 \\
  MuJoCo implicit & 331.7 & 677.4 \\
  MuJoCo RK4 & 1730.2 & 1861.8 \\
  Drake RK3 & 2531.6 & 5000.4 \\
}
\newcommand{\RotationFollowupRows}{%
Planar tilt & 0.732526 & 139.098 & 16.109 & 0.000122 & \(7.72\times10^{-5}\) \\
Curved tilt & 0.545388 & 1894.178 & 22.825 & 0.000122 & 0.000459 \\
Interior entry & 0.372769 & 157.350 & 105.379 & 0.000122 & 0.000176 \\
Rotating near miss & Clear & 29.091 & 0.017 & -- & -- \\
}
\newcommand{\HoleFollowupRows}{%
Counterbore passage & Clear & 939.851 & 934.138 & 0.380 & 31 \\
Counterbore shoulder & 0.744833 & 406.938 & 350.780 & 0.691 & 31 \\
Counterbore mouth & 0.244833 & 223.628 & 279.304 & 0.657 & 23 \\
Voron M5 spacer & Unresolved & 3042.649${}^{*}$ & 3782.196${}^{*}$ & 4635.623${}^{*}$ & 200000 \\
Voron M3 washer & 0.370605 & 2839.609${}^{*}$ & 1094.144 & 1035.655 & 46722 \\
Voron F695 bearing & 0.525604 & 4615.776${}^{*}$ & 5546.093 & 5222.019 & 171626 \\
}
\newcommand{\WitnessFollowupRows}{%
Trim-edge entry & 0.354330068 & hit & 0.354330068 & 0.00025 & 0.03108520603 \\
Narrow entry & 0.374969482 & hit & 0.374969479 & \(3.05\times10^{-5}\) & 0.0312499851 \\
Exact grazing & 0.375000000 & unresolved & 0.374999523 & -- & -- \\
Near miss & Clear & clear & 1.000000000 & -- & -- \\
}
\newcommand{\PerformanceRows}{%
  Cam array & B-rep & 0.724 & 0.0818 & 8.120 & 0 \\
   & Mesh (12) & 8.463 & 6.5374 & 0.504 & 188 \\
   & Mesh (24) & 11.422 & 1.2327 & 5.708 & 2 \\
   & Mesh (48) & 16.189 & 0.1210 & 7.946 & 31 \\
   & Mesh (96) & 21.662 & 0.0842 & 8.119 & 29 \\
  Ball train & B-rep & 1.077 & 1.5080 & 3.605 & 0 \\
   & Mesh (24) & 36.778 & 1.5080 & 3.605 & 0 \\
}
\newcommand{\ProductionRows}{%
  Wheel washer stack & 17/17 & 34.511 & 17.674 & \ensuremath{1.42\times10^{-13}} \\
  Wheel screw 1 & 5/17 & 1227.473 & 1.704 & \ensuremath{5.66\times10^{-5}} \\
  Wheel screw 2 & 4/17 & 1151.596 & 3.920 & \ensuremath{5.66\times10^{-5}} \\
  Lidar pre-seat & 0/17 & 886.659 & 18.289 & \ensuremath{5.64\times10^{-3}} \\
}

\title{\textbf{Certified Continuous Clearance of Trimmed STEP Surfaces under Fixed-Axis Rigid Motion}}
\author{NurbsCONTACT research manuscript}
\date{}

\begin{document}
\maketitle

\begin{abstract}
Continuous clearance queries must account for both the motion between poses and the regions removed by CAD trims. We certify positive clearance between nominal STEP boundaries under a prescribed relative fixed-axis rotation and affine translation. Outward-rounded directional bounds enclose the original curved control-point trajectories. For supported plane and cylinder faces under fixed relative orientation, radial and axial bounds expose the void inside a trimmed bore before parameter subdivision. The search returns a safe time prefix and, when an original-surface witness is available, an upper bound on first threshold entry. The circular-rim bounds reduce median runtime from 693.1 ms to 1.24 ms for a controlled eccentric bore. A counterbore reproduces this gain mechanism at the mouth and the internal shoulder. Tests on two production STEP sources show interface-dependent gains and work-cap limits. A conservative chord-based Chen OBB comparison uses the same analytic rotation references, and an independent audit checks trim-edge witnesses through near grazing. Separate implicit-dynamics experiments compare native and tessellated geometry in downstream models.
\end{abstract}

\section{Introduction}

An eccentric pin can approach a bore rim between two poses that are both safe. A convex enclosure of the entire face fills the bore, so a box or hull overlap may keep irrelevant times in the search. A continuous clearance query on the original STEP boundary must preserve the physical trim and enclose the path between the poses. This is particularly important when assembly verification specifies a positive margin.

We use two complementary exclusions. A directional bound encloses the scalar projection of each moving support point under exact fixed-axis motion. A radial and axial bound for supported face types uses the original trim to verify an empty disk and reject times in which the moving part stays in that void. The search keeps surviving time intervals through parameter subdivision. A feasible point pair on the trimmed surfaces supplies an independent upper bound on the first entry time.

Our contributions are an outward-rounded projection bound for the prescribed curved motion, a trim-dependent circular-void bound with a strict separation example against whole-face convex enclosures, and a first-entry certificate with an explicit unresolved outcome. We check these constructions on generated STEP cases with analytic references and on original component pairs from published assemblies. The downstream rigid-body tests compare STEP and tessellated geometry under shared HHT-$\alpha$ settings. Two analytically solvable contact problems also compare our simulation with native MuJoCo and Drake.

\section{Related work}

Static B-rep proximity methods retain geometric features that are approximated by faceted proxies. Crozet et~al. use smooth CAD boundaries and feature-specific routines to compute local minima for dynamics \cite{crozet2018}. Here the clearance query bounds a global minimum over the trimmed domains along a declared motion.

Chen et~al. use time-dependent inclusions for zero-distance collision of parametric surfaces, including rational B\'ezier patches \cite{chen2024}. Their Algorithms~2--3 propagate time intervals through subdivision and clip time using projections on AABB or OBB axes. These operations already establish time-interval inheritance and directional clipping. Their formulation gives linear control-point trajectories within a query, and their rigid-motion example uses time-linear pieces. Their termination returns the left endpoint of an earliest remaining candidate. Our query uses curved projection enclosures on the prescribed fixed-axis trajectory and reports a safe prefix at a clearance threshold with an optional feasible witness. Our circular-void predicate uses the native trim wires to tighten time exclusion on the supported faces.

Ferguson et~al. reduce curved rigid-body collision queries to conservative linear CCD with a separation margin \cite{ferguson2021}. This establishes prior use of positive margins for curved rigid motion. The present specialization concerns original trimmed STEP boundaries and their temporal bounds. The downstream comparisons use an AABB hierarchy for triangles \cite{larsen2000} and the HHT acceleration formulation \cite{hilber1977,chrono}.

\section{Certified continuous clearance}
\label{sec:continuous-clearance}

\subsection{Trimmed surfaces and the query contract}

Each face is \(S(D)\), where \(D\) is a closed parameter domain bounded by the
original STEP p-curves. A nominal connector closes a small UV endpoint gap
when needed, and the importer records its maximum enclosure. Boundary points
belong to \(D\). A parameter rectangle is discarded only after its separation
from every trim curve and its even crossing parity have both been certified.
Mixed rectangles remain in the hierarchy. Upper bounds use feasible points
on the original face or on its p-curves. Positive-weight NURBS spans are
subdivided into homogeneous B\'ezier nets with outward arithmetic. Native
elementary surfaces use analytic convex envelopes in their charts.

Let \(d(t)\) be the minimum distance between the two boundaries in the fixed
frame of body \(A\). Given a threshold \(\delta>0\) and interval \(I=[t_0,t_1]\),
we seek the first \(t\in I\) with \(d(t)\leq\delta\).
The lower time endpoint \(t_-\) certifies \(d(t)>\delta\) for all earlier
times in \(I\). A feasible original-surface point pair within \(\delta\)
at time \(t_+\) certifies an upper bound. An empty frontier with no witness
certifies the entire interval. If the budget is exhausted before a witness
is found, the remaining interval stays unresolved. A witnessed bracket can
also exceed the requested width, so the width is reported separately.

\subsection{Directional exclusion on the curved path}

Body \(B\) has an affine translation and rotation about a fixed unit axis
with angle \(\theta(t)=\theta_0+\omega t\) in this relative frame.
For a support vertex \(p_i\), Rodrigues' formula gives its projection
on a fixed direction \(n\) as
\[
f_i(t)=a_i+b_i t+c_i\cos(\theta(t))+d_i\sin(\theta(t)).
\]
Let \(H=\max_j n^\top a_j^A\) over the support vertices of the fixed
cell. Each time at which the two cell images are within \(\delta\) lies in
\(\bigcup_i\{t:f_i(t)\le H+\delta\|n\|\}\).

\begin{proposition}[Conservative directional time set]
Intersecting these sets over any collection of directions retains every
feasible time for the cell pair.
\end{proposition}
\begin{proof}
For feasible points \(x\) and \(y\) in the fixed and moving cells,
\(n^\top x\le H\) and \(n^\top(y-x)\le\|n\|\delta\).
Rigid motion preserves the convex weights of \(y\), so
\(\min_i f_i(t)\le n^\top y\). The inequality holds for every direction.
\end{proof}

An affine projection is clipped directly. For curved trajectories we
enclose the sine and cosine ranges with outward rounding and discard an
interval only if every lower support projection exceeds the upper
threshold. Ambiguous intervals survive the clipping budget. The retained
time set is passed to spatial descendants, as in time-dependent
inclusion methods \cite{chen2024}. The enclosure acts on the declared
curved trajectory throughout the interval.

\subsection{What the trim adds at a circular hole}

A planar face with a certified circular void has a positive lower radius
in a frame centered in the hole. To certify it, we enclose the minimum
distance from a proposed center to all trim curves in the UV chart.
The interior of that disk intersects no trim curve. Even parity at its
center therefore certifies that the whole disk lies outside the material
domain. A lower singular-value bound of the planar parameterization
translates the disk radius into physical units. Native cylindrical
faces supply radial and axial enclosures from their analytic charts.

\begin{proposition}[Radial and axial exclusion]
In a common cylindrical frame, let \(g_r\) and \(g_z\) be lower bounds
on the gaps between the two radial and axial coordinate intervals.
Then \(d\ge\sqrt{g_r^2+g_z^2}\). At fixed relative orientation, a
uniform \(g_r\le\delta\) restricts admissible axial separation to
\(\sqrt{\delta^2-g_r^2}\). If \(g_r>\delta\), the whole interval is safe
for the face pair.
\end{proposition}
\begin{proof}
The squared separation of any pair is the sum of its squared transverse
and axial separations. The reverse triangle inequality in the transverse
plane bounds the transverse separation below by the difference of the
two radii. Using the interval gaps proves the result. Under affine
translation and fixed orientation, the axial intervals admit affine
clipping with the stated allowance.
\end{proof}

\begin{proposition}[Strict separation from a whole-face convex bound]
Consider a planar annulus with inner radius \(R\) and outer radius
\(R_{\rm out}>R\). Let a disk of radius \(a\) remain coplanar during
the tested translation and let \(s(t)\) be its center offset. If
\(\sup_{t\in I}s(t)+a+\delta<R\), the radial test rejects all of \(I\),
while a single convex enclosure of the full annulus overlaps the disk
throughout \(I\).
\end{proposition}
\begin{proof}
The annulus has minimum radius \(R\), and every disk point has radius
at most \(s(t)+a\). Their radial gap exceeds \(\delta\). The outer circle
of the annulus spans the full disk of radius \(R_{\rm out}\) by convex
combinations. The moving disk lies inside that convex hull throughout.
\end{proof}

The example concerns a whole-face enclosure. Subdivision may eventually
recover the separation with convex cells. The void test can establish it
before those subdivisions, within its supported fixed-orientation regime.

\subsection{Preservation of the first-entry bracket}

Each queued task covers two spatial cells and a time set. Spatial and
temporal splits cover the parent domain, and the certified exclusions
remove only empty trim regions or times with distance above \(\delta\).
A valid witness time allows the search to ignore later events. A
budget-limited task keeps its earliest candidate time in the frontier.
Induction on these operations establishes the safe prefix and the
first-entry bracket whenever both endpoints are available. This argument
also applies to an interior event with safe trajectory endpoints.

\section{Continuous-query experiments}
\label{sec:continuous-experiments}

Six controlled fixtures cover planar entry, centered bore passage,
eccentric rim entry, rotary separation and entry, and a rotating spline
near miss. The first-entry time or no-event outcome follows independent
analytic geometry. In the archived matched-tolerance suite of 270 queries,
all 135 returned hit brackets enclose their reference time and all 135
no-event queries certify the full interval. The maximum hit width is
\(4.992\times 10^{-4}\) in normalized time. The tolerances are
\(2\times10^{-4}\) in space and \(5\times10^{-4}\) in time, with a
200,000-task cap. The componentwise chord-box ablation encloses each
rotating point's deviation from its linear chord by \(M h^2/8\), where
\(M\) bounds its second derivative and \(h\) is the interval width.
It runs in the shared search and evaluates the original path.

\begin{table}[!htbp]
\centering\small
\caption{Matched circular-rim ablation on a controlled fixture and original
production STEP assembly interfaces. Medians use nine query-only repeats.
Wheel first-entry brackets agree with independent OCCT cross-checks.}
\label{tab:rim-ablation}
\begin{tabular}{lrrr}
\toprule
Scenario & Directional (ms) & With radial and axial bounds (ms) & Ratio\\
\midrule
Controlled eccentric rim & 693.056 & 1.240 & 559\\
Centered bore passage & 16.957 & 0.403 & 42.1\\
Wheel washer stack & 23.887 & 7.605 & 3.14\\
Wheel eccentric bolt 1 & 905.886 & 151.779 & 5.97\\
Wheel eccentric bolt 2 & 925.171 & 330.183 & 2.80\\
Lidar nut pre-seat & 1647.266 & 1508.700 & 1.09\\
\bottomrule
\end{tabular}
\end{table}

In the controlled rim test, adding the circular bounds reduces visited
tasks from 33,728 to 25 and bracket width from \(4.992\times10^{-4}\)
to \(3.125\times10^{-5}\). Both brackets contain the analytic entry.
The production comparison uses the original STEP occurrences and placements
from the OpenAMRobot assemblies \cite{openamrobot}. The full lidar
insertion reaches the work cap and yields a wide witnessed bracket
\([0.386675,0.421701]\). The bracket overlaps the OCCT cross-check near
0.4001. Table~\ref{tab:rim-ablation} establishes a within-solver
exclusion benefit at the stated interfaces. Rechecking the archived
time outputs finds no disagreement in the analytic cases or with the
production reference intervals, including both wide lidar brackets.
The archive contains bracket endpoints and reference intervals.

\subsection{Common untrimmed patch comparison with TDIB-CCD}

We run the official Chen OBB implementation on the same untrimmed
biquadratic control nets and linear control-point trajectories as
our kernel. A planar and a curved moving graph first touch a fixed
plane at \(t=0.375\). A slower version of the curved graph remains
clear over \([0,1]\). Table~\ref{tab:chen-common} reports nine query-only
repetitions. Both methods use a numerical time target of
\(5\times10^{-4}\). The OBB width criterion also covers its parameter
cells, and the native query uses a distance tolerance of \(5\times10^{-4}\)
and a 25,000-task cap.

\begin{table}[!htbp]
\centering\small
\caption{Common zero-distance inputs. Chen's returned time is a candidate
at the left endpoint of a retained interval. Point residual uses the
returned cell centers. The native entry cases have a safe prefix with
no witnessed zero-distance upper bound at the cap.}
\label{tab:chen-common}
\begin{tabular}{lrrrr}
\toprule
Case & OBB candidate & OBB residual & Native safe prefix & Median ms (OBB / native)\\
\midrule
Planar entry & 0.375000 & \(4.92\times10^{-4}\) & 0.375000 & 218.0 / 1533.0\\
Curved entry & 0.375000 & \(6.91\times10^{-4}\) & 0.374832 & 132.4 / 1160.1\\
Curved clear & Clear & -- & 1.000000 & 2.93 / 1.08\\
\bottomrule
\end{tabular}
\end{table}

The native entry cases take longer under the reported stopping rules,
while both methods reject the clear case. These timings compare the
observed behavior of different query contracts. In a separate native query
at a margin of \(0.025\), both witnessed brackets enclose the analytic
first entry at \(0.35\). The planar and curved brackets have widths of
\(1.46\times10^{-4}\) and \(2.94\times10^{-4}\), respectively.

\ifdefined\RotationFollowupRows
\subsection{Conservative segmented-rotation comparison}
\label{sec:rotation-followup}

The follow-up uses a clearance of \(0.03125\) model units, a distance
tolerance of \(2\times10^{-4}\), a time target of \(5\times10^{-4}\) on
\([0,1]\), and a cap of 200,000 search tasks. Each configuration is
measured nine times after warm-up with alternating or rotating run order.
An 80-digit analytic evaluation provides the reference for the generated
geometries. The production pairs use the same independent adaptive OCCT
distance cross-check for all three methods, with a \(10^{-6}\) mm numerical
margin and a time resolution of \(2\times10^{-5}\).

For a segment of duration \(h\), the chord between the two exact endpoint
poses differs from the fixed-axis trajectory by at most \(M h^2/8\),
where \(M=\omega^2\max_i\|p_{i,\perp}\|\) bounds control-point acceleration.
The bound extends to each patch point by convexity. The segmented wrapper
splits until the chord bound fits the distance tolerance, then encloses the
rounding errors in endpoint position and velocity. It adds the resulting padding
and the requested clearance to the official Chen OBB projection test,
scaled by each axis norm. The wrapper visits segments in temporal order.
It subdivides retained time intervals before running the official
parameter search and verifies proposed rejections with outward projection
bounds on the chord nets. This check can refine the moving B\'ezier
net. Entering upper bounds require an original-path point pair.
The timing includes segmentation, projection verification, and witness
evaluation.

The fixtures share the same untrimmed biquadratic control nets. A moving
rectangle or curved graph tilts about the \(y\)-axis above a fixed plane.
At first entry the minimum vertical separation is
\(H-\sin(\omega t)\), with the closest moving point at \(x=1\)
and, for the graph, on its minimum-height centerline. One entry
occurs strictly between safe endpoints. The near miss stays above the
threshold by \(0.001\) model units at closest approach.

\begin{table}[!htbp]
\centering\small
\caption{Matched original-trajectory results. Every hit bracket contains
the analytic entry. Width is the largest across nine repeats.}
\label{tab:rotation-followup}
\begin{tabular}{lrrrrr}
\toprule
 & & \multicolumn{2}{c}{Median query (ms)} & \multicolumn{2}{c}{Bracket width}\\
\cmidrule(lr){3-4}\cmidrule(lr){5-6}
Fixture & Reference & Chen chords & Native & Chen chords & Native\\
\midrule
\RotationFollowupRows
\bottomrule
\end{tabular}
\end{table}

Both methods meet the time target on these fixtures. The curve forces
more segment and parameter checks in the chord method than a planar
fixture does. The native search certifies the rotating near miss with its
root directional bound. This comparison covers the padded and audited
rotation wrapper. The zero-distance linear comparison in
Table~\ref{tab:chen-common} measures the upstream implementation directly.

\subsection{Counterbore and a second production STEP source}
\label{sec:counterbore-followup}

The generated counterbore has mouth and throat radii of \(1.375\) mm
and \(0.8125\) mm, with an internal shoulder at their junction and a
pin of radius \(0.375\) mm. Different lateral offsets lead to a clear passage, a first
entry at the shoulder, or a first entry at the mouth. For each circular
rim, the reference pairs the leading pin cap with that rim and combines
the radial and axial gaps. At the two entry offsets the radial gap at the
active rim is \(0.0234375\) mm.

The production replication uses the Voron Trident Front Idler B
assembly \cite{vorontrident}. The extractor keeps original AP214
geometry and occurrence transforms within the idler subtree. We test the
M5 spacer and an F695 bearing against their original M5 screw, and the
M3 washer against its screw. Each approach follows the screw axis with a
transverse displacement of \(0.10\) mm. The three variants share the
same trim handling, search, clearance and reference. They use a swept
box, directional time bounds, or directional time bounds with radial
and axial exclusion.

\begin{table}[!htbp]
\centering\small
\caption{Three-way ablation. The reference column gives the analytic
entry for the counterbore and the lower endpoint of the OCCT bracket for
the Voron interfaces. A star marks a budget-limited result or a bracket
wider than the time target.}
\label{tab:counterbore-followup}
\begin{tabular}{lrrrrr}
\toprule
 & & \multicolumn{3}{c}{Median query (ms)} & \\
\cmidrule(lr){3-5}
Fixture & Reference & Swept box & Directional & With radial/axial & Tasks with bound\\
\midrule
\HoleFollowupRows
\bottomrule
\end{tabular}
\end{table}

The radial and axial filter leaves 23--31 tasks for each counterbore
query, compared with 13,754--44,395 for directional clipping alone.
All three variants agree with the analytic references. For the Voron
washer and bearing, both directional variants meet the time target,
while the swept-box variant reaches the cap with wider witnessed
brackets. The radial and axial filter reduces task counts modestly.
All three methods exhaust their budget on the M5 spacer. The best
certified prefix reaches \(0.523320\), and the independent numerical
reference reaches \(0.975586\) before its own cap with no upper witness.

\subsection{Independent trim-edge certificate audit}
\label{sec:edge-audit}

For a plane with a unit circular trim and a parallel disk of radius
\(0.25\), the disk stays inside the hole in projection. At closest distance
the plate point lies in the interior of the inner trim edge at
\((1,0,0)\). The analytic
distance is \(\sqrt{g^2+(0.375-t)^2}\), with \(g\) the in-plane gap.
Offsets place \(g\) well inside the threshold, \(2^{-26}\) below it,
exactly on it, or \(2^{-23}\) above it. The narrow entry interval is
about \(6.10\times10^{-5}\) wide and both trajectory endpoints are safe
in every case.

The independent verifier compares the returned prefixes and witnessed
times with the analytic reference and reconstructs the residual from
the exported Cartesian coordinate intervals. It checks those intervals
against the two plane equations and the disk and trim radii. The
threshold equality at \(t=0.375\) is an isolated contact. A prefix with no
feasible upper witness leaves its remaining interval unresolved.

\begin{table}[!htbp]
\centering\small
\caption{Independent trim-edge audit. The point-pair residual is an
upper bound recomputed from the archived coordinate intervals.}
\label{tab:edge-audit}
\begin{tabular}{lrrrrr}
\toprule
Case & Reference & Outcome & Safe prefix & Width & Residual\\
\midrule
\WitnessFollowupRows
\bottomrule
\end{tabular}
\end{table}

Both crossing cases return edge witnesses and contain their analytic
entry times. The exact graze remains unresolved, with its prefix ending
\(4.77\times10^{-7}\) before the reference time. The near miss certifies
the entire interval at a minimum margin of \(2^{-23}\) mm.

\fi


\section{Application to implicit multibody contact}

\subsection{Reduced coordinates}

The position of a local point \(x\) on body \(i\) is prescribed by its one-degree-of-freedom joint. For a slider it is \(o_i+e_i q_i+x\). For a hinge it is \(o_i+R_z(q_i)x\). The coordinate mass is the body mass or its moment of inertia about the hinge. Fixed bodies have constant placements. These kinematics yield a constant positive diagonal matrix \(M\). Couplings enter the generalized force through linear springs, dampers, drive targets, and contact. In the cam array, six lateral stations have disjoint \(x\) intervals, and each cam interacts with its own follower. The ball train includes all sphere pairs and all sphere-to-stop pairs.

For each eligible pair let \(d\) be the boundary distance and let \(p_A,p_B\) be a closest pair with \(d>0\). With \(n=(p_B-p_A)/\norm{p_B-p_A}\), the gap Jacobian projects each body's point velocity onto \(n\) with the appropriate sign. The velocity derivatives for a slider and hinge are \(e_i\) and \(e_z\times(p_i-o_i)\), respectively. A linear compliance layer of thickness \(s\) generates the force
\begin{equation}
    F_c=J^\top\lambda,\qquad
    J=\frac{\partial d}{\partial \q},\qquad
    \lambda=k[s-d]_+.
    \label{eq:compliance}
\end{equation}
This choice exposes geometry error through both the activation threshold and the load. The analytic and triangulated backends share \(s\), \(k\), and the same generalized force assembly.

\subsection{Analytic B-rep fast path}

We extract the analytic frames and extents from STEP. The fixture loader checks the finite rectangle and the full round solids against their face structure and source geometry. For two spheres, the nearest points lie on their center line whenever they are positively separated. For a sphere or cylinder facing a bounded planar patch \(Q\), first obtain the smallest signed support distance \(d_\pi\) to the infinite supporting plane \(\pi\). For a cylinder with unit axis \(a\), radius \(r\), half-length \(\ell\), and center \(c\), it is
\begin{equation}
d_\pi=n^\top(c-c_Q)-r\norm{n\times a}-\ell|n^\top a|.
\end{equation}
The two radial coordinates and axial coordinate identify a minimizing native point \(p_B\). When \(n^\top a=0\) we select its axial midpoint.

\begin{proposition}[Finite-face support]
If \(d_\pi>0\) and the orthogonal projection of \(p_B\) lies in \(Q\), then the global boundary distance to \(Q\) equals \(d_\pi\). The projected point and \(p_B\) form a feasible witness.
\end{proposition}
\begin{proof}
For every \(q\in Q\) and point \(b\) on the round body, the component of \(b-q\) along \(n\) is at least \(d_\pi\). At \(p_B\) and its valid projection the segment is parallel to \(n\) and has length \(d_\pi\). The lower bound and the witness agree.
\end{proof}

The membership check uses the finite source rectangle. When the condition fails, a generic trimmed B-rep distance query uses the original p-curves and bounded surface evaluations \cite{chen2024,crozet2018}. Its lower and upper distance bounds can report an unresolved budget. The dynamic evaluation reports that condition as a failure to obtain a witness within tolerance.

\subsection{Triangle AABB baseline}

We triangulate the same STEP spheres and cylinders with uniform azimuthal sampling. A sphere uses half as many latitudinal intervals as azimuthal segments. Each finite rectangle uses two triangles. The general trimmed-query experiment instead uses Open CASCADE tessellation with a requested linear deflection. A median-split local AABB tree encloses the triangles with two primitives per leaf. During a query the moving node bounds are transformed into the fixed body's frame. Best-first traversal uses box separation as a lower bound and a feasible triangle pair as an upper bound. The leaf calculation covers vertex-to-face, edge-to-edge, and segment--triangle intersection cases. A previous triangle pair can seed the upper bound. Both coherence settings use the same traversal and complete the same distance problem for their triangulation.

\subsection{Implicit integration and acceptance}

The unknown at each step is \(\acc_{k+1}\). At step size \(h\), we apply the Newmark updates
\begin{align}
    \q_{k+1}&=\q_k+h\vct_k+h^2\bigl[(1/2-\beta)\acc_k+\beta\acc_{k+1}\bigr],\\
    \vct_{k+1}&=\vct_k+h\bigl[(1-\gamma)\acc_k+\gamma\acc_{k+1}\bigr],\\
    \mathbf R&=M\acc_{k+1}-(1+\alpha)\force(\q_{k+1},\vct_{k+1},t_{k+1})
                    +\alpha\force(\q_k,\vct_k,t_k).
\end{align}
We use \(\alpha=-0.1\), \(\beta=(1-\alpha)^2/4\), and \(\gamma=1/2-\alpha\). In the acceleration solve, the symmetric tangent is
\begin{equation}
M+(1+\alpha)\bigl[\beta h^2 K+\gamma h C\bigr],
\end{equation}
where \(K\) combines spring derivatives and the active contact approximation \(kJ^\top J\), and \(C\) contains the viscous couplings. The full nonlinear force is recomputed in the residual. A backtracking line search accepts a residual decrease. A converged residual satisfies an absolute tolerance of \(10^{-8}\) plus \(10^{-8}\) times the larger force infinity norm at the two endpoints. If the iteration or clearance check fails, the step is retried as two half steps. A failed subdivision sequence rolls back the outer step.

\begin{proposition}[Positive clearance of the reconstructed step]
For each contact pair, let \(\underline d_0,\underline d_1\) enclose the exact distance from below at the two endpoints. Let \(L_A,L_B\) be upper bounds on point travel under linear interpolation of the reduced coordinates. Then the interval-wide distance is bounded below by
\begin{equation}
    \underline d_{\mathrm{step}}
      =\min(\underline d_0,\underline d_1)-\tfrac12(L_A+L_B).
    \label{eq:guard}
\end{equation}
\end{proposition}
\begin{proof}
Boundary distance is Lipschitz with respect to the sum of the two Hausdorff displacements. The path from the nearest endpoint has length at most \((L_A+L_B)/2\) because the slider speed and angular speed bound are constant under the declared interpolation. Apply the endpoint lower bound at that endpoint.
\end{proof}

For a slider we bound travel by the position change, and for a hinge by the angular change times an enclosing distance to the pivot. The analytic endpoint bounds and travel bounds use outward interval arithmetic. We require \(\underline d_{\mathrm{step}}>0.5\) mm for every eligible pair. Both pipelines use this CAD guard as a common acceptance cost. A cache keyed by the exact previous coordinate vector reuses endpoint bounds. Cache hits change only the number of evaluations. The guarantee concerns the linear reconstruction between accepted states on the nominal CAD fixtures.

\section{Rigid multibody numerical benchmarks}

\begin{figure}[t]
\centering
\includegraphics[width=\linewidth]{model-scenes.pdf}
\caption{STEP-derived dynamic fixtures rendered from accepted HHT states. The driven cam array couples six hinges and six compliant sliders. The ball train couples twelve sliders through sphere contacts and the two stops.}
\label{fig:models}
\end{figure}

\subsection{Protocol and reference problems}

We run Release builds on a Windows x64 workstation with an AMD Ryzen 9 5900X, MSVC 19.44, and Open CASCADE 7.9.3. The dynamic model uses meters, seconds, and SI masses. The mesh and native variants receive the same source STEP files, HHT parameters, joint graph, force law, and acceptance tolerances. For the long compared simulations, timings are medians of three successive compute-only repetitions per configuration. Model loading, tessellation, and writing trajectories fall outside the timer. Integration, nonlinear forces, geometry queries, and the common clearance guard fall inside. The raw records retain repeat-level spread and the work counts. Ball-train comparisons use the faster full-length mesh coherence setting. A separate ablation uses five repetitions for the guard cache.

The linear oscillator has analytic solution \(q(t)=\cos(10t)\). It checks accuracy independently of the geometry. The single-ball case starts with speed \(0.45\) m/s toward a fixed floor, a \(0.018\) m stand-off beyond the compliance layer, unit mass, and no gravity. After entry, the compression is a half sine at frequency \(\sqrt{k/m}\). The contact force is \(k\) times that compression. We also sample the eccentric cam through a full turn. Its closed-form clearance is \(0.143-0.105-0.024\sin\theta\) m.

The cam array runs for 120 s at a 5 ms nominal step. Adjacent hinge coordinates and adjacent sliders are elastically coupled. The ball train runs for 60 s at the same step under a two-frequency input force on the first ball. Its positive-gap contacts transmit motion through the train. We also integrate both native cases at 2.5 and 1.25 ms to assess temporal sensitivity. The smallest step supplies the numerical trajectory reference. Two 300 s native runs check sustained operation. The contact layer thickness is 12 mm and \(k=5000\) N/m for all trajectories.

\subsection{Calibration and faceting error}

\begin{figure}[t]
\centering
\includegraphics[width=\linewidth]{calibration.pdf}
\caption{Calibration separates time integration and contact geometry. \textbf{a}, Analytic oscillator error converges at second order. \textbf{b}, The computed contact force follows the half-sine solution. \textbf{c}, The maximum geometric clearance error drops with mesh resolution. The native cam error is below the plotted range.}
\label{fig:calibration}
\end{figure}

The observed oscillator orders were \OscillatorOrders{}. The native impact displacement RMS error was \DropError{} mm at 1 ms and \DropFineError{} mm at 0.5 ms. Figure~\ref{fig:calibration} compares the force response with the analytic curve. The paired 1 ms pole-aligned sphere mesh produced the same trajectory to the numerical tolerance because it contains the required supporting vertex.

On the rotating cam, inscribed polygons overestimate the plane clearance by up to \(r[1-\cos(\pi/N)]\), where \(N\) is the number of azimuthal facets. The nearest available vertex can be offset from the support direction by at most \(\pi/N\). The measurements are consistent with the resulting quadratic decay in angular facet size. Meshes from 12 to 96 facets reduce the observed maximum error from \CamCoarseGapError{} to \CamFineGapError{} mm over 360 poses per configuration. The analytic B-rep clearance matches the reference to roundoff. The generic B-rep search reached its 30,000-task cap before closing the requested \(2\,\mu\mathrm m\) interval on all twelve tested cam poses. This motivates the certified finite-face specialization.

\subsection{Multibody cost and accuracy}

\begin{table}[t]
\centering\small
\caption{Measured multibody outcomes. Errors are translational RMS differences against finer-step native runs. Time is the median of three repetitions. The guard column is the smallest interval lower bound across those repetitions.}
\label{tab:performance}
\begin{tabular}{llrrrr}
\toprule
Scene & Backend & Runtime (s) & RMS (mm) & Guard (mm) & Subdivisions \\
\midrule
\PerformanceRows
\bottomrule
\end{tabular}
\end{table}

\begin{figure}[t]
\centering
\includegraphics[width=\linewidth]{cost-accuracy.pdf}
\caption{Cost and accuracy under a shared HHT solver. \textbf{a}, Ball-train medians and repeat ranges at matched trajectory accuracy. \textbf{b}, Cam-array time against follower-position RMS error relative to the 1.25 ms native trajectory. Mesh points are annotated by azimuthal facet count.}
\label{fig:cost}
\end{figure}

The 60 s ball-train results give a comparison within a \(1\,\mu\mathrm m\) trajectory tolerance. The sphere-to-sphere normals are collinear with the slider axis, which the tested mesh samples directly. The largest difference between the native trajectory and the faster mesh configuration was \TrainMatchError{} mm over all sliders and recorded outer-step times. The measured median ratio was \TrainGain{} times. Both configurations use the same positive clearance check. The two mesh variants produced identical trajectories. The short unseeded traversal performed fewer primitive tests and had a slightly lower median latency. Timing varied across batches, so the main comparison retains the fastest measured full-length baseline.

For the cam array, geometric error propagates to the follower even under identical HHT steps. Table~\ref{tab:performance} and Figure~\ref{fig:cost} show that denser meshes reduce the trajectory discrepancy at higher computational cost. At 96 facets, the follower RMS difference relative to the finer-step native reference was \FineCamMeshError{} mm. The matched nominal-step native trajectory had an RMS difference of \CamCadTimeError{} mm. The native run was \CamFineGain{} times faster than this mesh configuration. The finest mesh had a smaller maximum follower discrepancy despite its larger RMS difference. These are errors relative to the stated numerical reference, with analytic static clearance evidence in Figure~\ref{fig:calibration}.

Halving the native nominal step reduced the cam array RMS error from \CamCadTimeError{} to \CamTemporalFine{} mm. The same refinement reduced the first-five-second ball-train RMS error from \TrainEarlyCoarse{} to \TrainEarlyFine{} mm. Over the full 60 s the corresponding ball-train discrepancies were \TrainLongCoarse{} and \TrainLongFine{} mm against the 1.25 ms trajectory. The long-window error is nonmonotonic under time-step refinement. Backend agreement at a common time step, analytic oscillator convergence, and long-time trajectory accuracy therefore measure different properties.

Reusing endpoint guards leaves every tested cam state unchanged. Over the 18 s ablation it reduces endpoint evaluations from \GuardOff{} to \GuardOn{}, and the median time falls by a factor of \GuardGain{}. The cache is shared by the two geometric pipelines. Its savings are therefore independent of which narrow-phase backend generated the force.

\subsection{Long-time behavior and a mesh-favorable case}

\begin{figure}[t]
\centering
\includegraphics[width=\linewidth]{long-run.pdf}
\caption{Two 300 s native HHT runs. The line marks the mean of the recorded pose-level metric per one-second bin, and the band spans the minimum and maximum within the bin. The independent interval-wide guard minima are reported in the text.}
\label{fig:long}
\end{figure}

Each 300 s native run completed 60,000 nominal advances. Minimum interval-wide clearances were \StressCamGuard{} mm for the cam array and \StressTrainGuard{} mm for the train. Figure~\ref{fig:long} shows contact load and clearance with one-second ranges. These observations establish sustained completion and positive clearance for the evaluated configurations. The runs exercise repeated forcing and contact rather than a linear stability theorem for arbitrary nonlinear systems.

We separately compare static clearance on a STEP pin and trimmed circular bore at 25 poses. The original B-rep retains the radial bore boundary and encloses the analytic distance at each pose. With the tested finest tessellation the triangle BVH achieved a maximum point estimate error of \TrimFinestError{} mm against that reference, and a median time of \TrimFinestMesh{} ms. The general B-rep median was \TrimCadTime{} ms. The dataset demonstrates two regimes. Analytic native support can improve both cost and accuracy on qualified pairs. A tessellated hierarchy can be much faster on a general trimmed shape at the tested accuracy.

\subsection{Additional static STEP samples}

\begin{table}[t]
\centering\small
\caption{Static clearance on original production STEP occurrences. The timings are medians across 17 shared poses per assembly pair. The mesh uses requested deflection 0.0125 mm. Its error column reports maximum pointwise distance discrepancy against OCCT boundary extrema. The native median includes any queries that reached the task budget.}
\label{tab:production}
\begin{tabular}{lrrrr}
\toprule
Assembly motion & CAD resolved & CAD (ms) & Mesh (ms) & Mesh error (mm)\\
\midrule
\ProductionRows
\bottomrule
\end{tabular}
\end{table}

We additionally sample 17 poses per production motion and directly compare the native B-rep search and the mesh BVH. Both use original placements and boundary distances in millimeters. The B-rep search requests a 0.0002 mm interval width with a cap of 30,000 tasks. Each of the mesh deflections 0.20, 0.05, and 0.0125 mm is tessellated independently from the source shape. The numerical reference is OCCT boundary extrema at the same pose. The CAD intervals contain the numerical reference at \ProductionOracleAgreement{} of \ProductionOraclePoses{} poses with a \(10^{-6}\) mm comparison margin. The finest mesh returned a point estimate at every pose and had the smaller median query time for each pair, with the largest observed discrepancy on the lidar support (Table~\ref{tab:production}). Section~\ref{sec:continuous-experiments} reports the separate trajectory-wide results.

\ifdefined\EngineBenchmarkRows
\subsection{Open-source rigid-body simulators}
\label{sec:engine-comparison}

We compare the STEP/HHT simulation with native MuJoCo~\cite{mujoco}
and Drake~\cite{drake}. The two controlled problems are a sphere
approaching a plane and a collision of two translating spheres.
Each mobile body has mass 1 kg. The approaching sphere has initial
speed 0.45 m/s. Both problems start 18 mm outside a 12 mm contact
layer and use an undamped normal stiffness of 5000 N/m, with
gravity and friction set to zero. The analytical reference uses
the one-contact half sine with reduced mass 1 or 0.5 kg. The
two-sphere case also conserves center-of-mass velocity.

The layer is represented by offset collision surfaces. The plane
is raised by the layer thickness in the first problem. In the
second problem each collision radius grows by half that thickness.
The three simulators use these offsets in their own contact
representations. Our directional STEP probes evaluate the offset
surfaces, and the trajectory-wide clearance tests are reported
separately in Section~\ref{sec:continuous-experiments}.
MuJoCo \EngineMuJoCoVersion{} uses its native one-dimensional
contact and direct reference stiffness scaled by the contact
inverse mass. Drake \EngineDrakeVersion{} uses native continuous
point contact with zero Hunt--Crossley dissipation and geometry
stiffnesses of 10,000 N/m in series.
The maximum force-law residual on the computed states is
\EngineForceResidual{} N across the tested steps.

\begin{table}[!htbp]
\centering\small
\caption{Native rigid-body contact at a 1 ms nominal step. Position
errors use all recorded coordinates and the analytic trajectory.
Force errors use the same sample times and the analytic pulse.}
\label{tab:engine-error}
\begin{tabular}{lrrrr}
\toprule
 & \multicolumn{2}{c}{Position RMSE (\(\mu\)m)} &
 \multicolumn{2}{c}{Force RMSE (N)} \\
\cmidrule(lr){2-3}\cmidrule(lr){4-5}
Method & Sphere--plane & Two spheres & Sphere--plane & Two spheres \\
\midrule
\EngineBenchmarkRows
\bottomrule
\end{tabular}
\end{table}

\begin{figure}[!htbp]
\centering
\includegraphics[width=\linewidth]{\EngineFigurePath}
\caption{Rigid-body comparison with native MuJoCo and Drake.
All trajectories use the same isolated normal force law.
Position errors are evaluated against the analytical solution.
The force panel shows the sphere--plane pulse at 1 ms.}
\label{fig:engine-comparison}
\end{figure}

Table~\ref{tab:engine-error} and Figure~\ref{fig:engine-comparison}
show the effect of the integration method on the same force pulse.
At 1 ms, MuJoCo RK4 and Drake RK3 have smaller position errors
than HHT in both examples. The archived runs include all three
steps, seven integration timings per configuration, and every
scored state. The wider CAD experiments evaluate different costs,
including trim classification and interval exclusion.

\ifdefined\EngineTimingRows
\begin{table}[!htbp]
\centering\small
\caption{Median integration time for 160 steps at 1 ms in the
same Ubuntu WSL environment. Model setup and trajectory export
are outside the timed region.}
\label{tab:engine-time}
\begin{tabular}{lrr}
\toprule
Method & Sphere--plane (\(\mu\)s) & Two spheres (\(\mu\)s) \\
\midrule
\EngineTimingRows
\bottomrule
\end{tabular}
\end{table}
\else
The STEP/HHT timings use Windows, and both engine timings use Ubuntu
WSL. The integration timings are archived with their environments.
\fi

\fi

\section{Scope and conclusion}

Fixed-axis projection enclosures and the trim-derived circular void
retain distinct information about the original STEP boundaries. Both
preserve coverage of the first-entry search under their stated motion
and face assumptions. The strongest measured gains occur when the
void removes most of a coarse candidate interval, including both
entry rims of the counterbore. The Voron spacer and full lidar insertion
show where the work cap limits the available certificate. The trim-edge
audit confirms threshold crossings and shows the finite safe prefix
returned at isolated grazing.

The trajectory certificate concerns a specified relative motion and
nominal closed trim domains. The HHT application uses its own accepted
configuration reconstruction and endpoint travel bounds. Its numerical
comparisons quantify trajectory accuracy and computational cost for
rigid bodies under the stated contact and integration settings.

\paragraph{Reproducibility.}
The MuJoCo and Drake scripts are in \path{benchmarks/engines/}. The trajectory and timing archive is in \path{results/engine-comparison/}.
The paired eccentric-rim measurements are in \path{benchmarks/axial-results.csv}. The follow-up protocol and raw results are in \path{benchmarks/clearance-followup/}. The production STEP input hashes and motion definitions are in \path{benchmarks/assemblies/}. The \path{nc_hht} target generates and reads the dynamics STEP fixtures. \path{benchmarks/hht/run.py} runs the dynamics suite, \path{analyze.py} validates trajectories and aggregates results, and \path{render.py} and \path{plots.py} generate videos and figures. \path{results/hht/protocol.json} specifies the run matrix and the results contain the multi-view animation index.

\bibliographystyle{unsrtnat}
\bibliography{manuscript}
\end{document}
